
\subsubsection{2.1 Data Curation for Pretraining}

DCLM \citep{Li2024DCLM} establishes that fastText quality filtering of CommonCrawl yields substantial performance improvements (64\% MMLU at 7B scale), but does not evaluate demographic or fairness implications of filtering decisions. FineWeb \citep{Penedo2024FineWeb} achieves competitive performance via C4-like filtering plus deduplication on 15T tokens, also without fairness evaluation. DoReMi \citep{Xie2023DoReMi} introduces domain reweighting via a proxy model, demonstrating that domain composition is a lever for downstream performance, but again measuring only performance outcomes. Dolma \citep{Soldaini2024Dolma} provides modular filtering infrastructure and documents demographic limitations (primarily English, Western-centric), without quantifying how specific filtering decisions transform demographic-occupation associations.

\subsubsection{2.2 Fairness in Language Models}

Bender et al. (2021) argue qualitatively that large-scale web corpora encode demographic harms through statistical patterns. BBQ \citep{Parrish2022BBQ} measures social biases in question answering; WinoBias \citep{Zhao2018WinoBias} targets gender-stereotyped coreference resolution; StereoSet \citep{Nadeem2021StereoSet} measures stereotype endorsement. These benchmarks characterize model output behavior but do not establish connections to specific pretraining data curation decisions.

\subsubsection{2.3 Data Documentation}

Data documentation frameworks (Gebru et al., 2018; Bender and Friedman, 2018) advocate for structured disclosure of dataset properties but describe static corpus characteristics rather than measuring how curation operations transform those properties.

\subsubsection{2.4 Positioning}

This work extends the data curation evaluation framework to include demographic association metrics at corpus level, and provides quantitative measurement rather than qualitative documentation of how curation operations alter demographic structure.

